\chapter{梯度、海森矩阵及数值实现细节}

\section{导数推导的意义}
非线性优化性能高度依赖导数信息质量。若梯度与海森矩阵计算准确，则牛顿法具有更快的局部收敛速度；法方程方法也可借助雅可比构造稳定的近似二阶信息。课程代码中对测距、方位角、俯仰角函数分别实现了一阶与二阶导数，本章给出其数学对应关系。

\section{测距函数导数}
记
\[
r=\sqrt{\Delta x^2+\Delta y^2+\Delta z^2}.
\]
则测距函数梯度为
\begin{equation}
\nabla \rho=
\left[
\frac{\Delta x}{r},
\frac{\Delta y}{r},
\frac{\Delta z}{r}
\right]^{\T}.
\label{eq:dist_grad}
\end{equation}
海森矩阵为
\begin{equation}
\nabla^2\rho=
\frac{1}{r}\bm I_3-\frac{1}{r^3}
\begin{bmatrix}
\Delta x^2 & \Delta x\Delta y & \Delta x\Delta z\\
\Delta y\Delta x & \Delta y^2 & \Delta y\Delta z\\
\Delta z\Delta x & \Delta z\Delta y & \Delta z^2
\end{bmatrix}.
\label{eq:dist_hess}
\end{equation}
该结果与代码 \texttt{distgrad()}、\texttt{disthess()} 完全一致。

\section{角度函数导数}
\subsection{方位角梯度}
定义 $s=\Delta x^2+\Delta y^2$，则
\begin{equation}
\frac{\partial A}{\partial x}=\frac{\Delta y}{s},\quad
\frac{\partial A}{\partial y}=-\frac{\Delta x}{s},\quad
\frac{\partial A}{\partial z}=0.
\end{equation}

\subsection{俯仰角梯度}
令 $r_{xy}=\sqrt{\Delta x^2+\Delta y^2}$，$r=\sqrt{\Delta x^2+\Delta y^2+\Delta z^2}$，则
\begin{equation}
\frac{\partial h}{\partial x}=-\frac{\Delta x\Delta z}{r_{xy}r^2},\quad
\frac{\partial h}{\partial y}=-\frac{\Delta y\Delta z}{r_{xy}r^2},\quad
\frac{\partial h}{\partial z}=\frac{r_{xy}}{r^2}.
\end{equation}
代码中在 $r$ 过小情况下设置异常标记，以避免除零风险。

\subsection{二阶导数实现}
角度函数二阶导数表达较长，程序中通过显式公式给出，关键考虑包括：
\begin{itemize}
  \item 利用对称性减少重复计算；
  \item 将公共项提取为中间变量，提高效率；
  \item 在近奇异几何（目标接近测站）时触发保护逻辑。
\end{itemize}

\section{目标函数梯度与海森矩阵}
目标函数可写为
\[
F(\bm p)=\sum_i w_i r_i(\bm p)^2,\quad w_i=\sigma_i^{-2}.
\]
因此梯度为
\begin{equation}
\nabla F(\bm p)= -2\sum_i w_i r_i(\bm p)\nabla f_i(\bm p),
\end{equation}
其中 $f_i$ 表示理论观测函数。海森矩阵为
\begin{equation}
\nabla^2 F(\bm p)=
2\sum_i w_i\left[\nabla f_i\nabla f_i^{\T}-r_i\nabla^2 f_i\right].
\label{eq:hessian_full}
\end{equation}
式 \eqref{eq:hessian_full} 与程序 \texttt{hessfun()} 的实现一一对应。

\section{线性方程求解与矩阵逆}
代码中对 $3\times3$ 线性方程组采用显式代数公式求解，并给出 $3\times3$ 矩阵伴随法求逆。该策略在小规模问题中计算量低、实现直接，但需要注意：
\begin{enumerate}
  \item 若矩阵接近奇异，显式公式数值稳定性下降；
  \item 实际工程建议增加条件数监测；
  \item 在高维扩展问题中应改用 Cholesky、QR 或 SVD 分解。
\end{enumerate}

\section{程序模块与数学模块对应关系}
\begin{table}[H]
\centering
\caption{主要函数与数学角色对应}
\begin{tabular}{ll}
\toprule
函数名 & 数学含义 \\
\midrule
\texttt{dist()}, \texttt{dire()} & 观测函数 $f_i(\bm p)$ \\
\texttt{distgrad()}, \texttt{diregrad()} & 雅可比（梯度）分量 \\
\texttt{disthess()}, \texttt{direhess()} & 二阶导数分量 \\
\texttt{objfun()} & 目标函数 $F(\bm p)$ \\
\texttt{gradfun()} & 梯度 $\nabla F(\bm p)$ \\
\texttt{hessfun()} & 海森矩阵 $\nabla^2F(\bm p)$ \\
\texttt{AXB()} & 线性系统求解器 \\
\texttt{AInverse3()} & 小维矩阵逆 \\
\bottomrule
\end{tabular}
\end{table}

\section{本章小结}
本章建立了导数计算与程序实现之间的可核查映射。通过解析导数而非数值差分，算法在效率和精度上均具优势，也为后续收敛性分析和不确定性评估提供了可靠基础。
